\cotefile{84}
\scanpage{113}

% manuscript lines / position: 1–2
Je dis que cette formule implique que le diagramme similaire au précédent

% manuscript lines / position: 3–6
\diagram{\begin{tikzcd}[ampersand replacement=\&,column sep=large]
 A_1\otimes_k A_2 \arrow[r] \arrow[d,"{\deleted{\unsure{N_1}}\,n_1\otimes_k\mathrm{id}_{A_2}}"'] \arrow[d,phantom,"{=N_{12}}" {xshift=2em}] \& A \arrow[d,"N_2",no head]\\
 A_2 \arrow[r,shift left=.35ex,no head] \arrow[r,shift right=.35ex] \& A_2
\end{tikzcd}}

% manuscript lines / position: 4–7, left of diagram, arrow from left label
\margin{déduite de la forme $n_1$ par chgt de base $k\longrightarrow A_2$}

% manuscript lines / position: 8–10
C'est en tout cas vrai pour les éléments de \deleted{la forme} $A_1\otimes_k A_2$ de la forme $x_1\otimes x_2$, car les deux mbres sont alors

% manuscript lines / position: 11
\[
 t_1(x_1)\cdot\deleted{\unsure{\alpha(x_2)}}x_2^2
\]

% manuscript lines / position: 12–15
par quadraticité de $N_{1,2}$, $N_2$. Il reste à voir qu'on a commutativité pour les \deleted{formes} bil. associées, or pour $N_{12}$ c'est la forme déduite de

% manuscript lines / position: 16
\[
 \deleted{(x,y)\longrightarrow}\quad \varphi_{n_1}(x,y)=t_1(x\bar y)\qquad\text{sur } A_1\times A_1
\]

% manuscript lines / position: 17–18
par chgt. de base $k\longrightarrow A_1$, or on a si \inserted{$x'=\alpha_1(x),\,y'=\alpha_1(y)$}

% manuscript lines / position: 18–20
\[
 \left\{\begin{aligned}
 \deleted{\unsure{x'}}\,\alpha_1(\bar y)&=\sigma_2(\alpha_1(y))\\
 \underset{\text{dans }A_2}{t_1(x\bar y)}&=T_2(\deleted{\unsure{\alpha|\alpha_1}}(x\bar y))\; \unsure{.}
 \end{aligned}\right.
\]

% manuscript lines / position: 21–22
\[
 \text{donc}\qquad \varphi_{n_1}(x,y)=\underset{\substack{\unsure{\mathord{''}|=}\\ \varphi_{N_2}(x',y')}}{T_2(\unsure{\cdot}x'\sigma_2(y'))}\qquad \text{O.K.}
\]

% manuscript lines / position: bottom left
% archive sheet number: 112
\scanpage{114}

% manuscript lines / position: 1–3
\uline{NB} la condition sur $N_2$ implique celle sur $T_2$, en passant aux formes bil. ass., i.e. en l'appliquant pour $x$ \uline{et} pour $1+x$ $(x\in A_1)$.

% manuscript lines / position: 4
\textcircled{$3^\circ$} On voit la condition

% manuscript lines / position: 5
\[
 \boxed{\sigma_1\sigma_2\sigma_3=1}
\]

% manuscript lines / position: 4–9, left, oblique
% slanted
\margin{Cette condition est automatiquement vérifiée sur la sous-algèbre de $A$ engendrée par les $A_i$, car elle est vraie sur \unsure{chaque|les} $A_i$ (lequel est stable par les $\sigma_i$).}

% manuscript lines / position: 6–8
implique que les involutions $\sigma_i$ commutent entre elles, car le produit de 2 d'entre elles est une involution

% manuscript lines / position: 8–10
\[
 \boxed{\begin{aligned}\sigma_1\sigma_2&=\sigma_3\\
 \sigma_2\sigma_3&=\sigma_1\\
 \unsure{\cdot}\sigma_3\sigma_1&=\sigma_2\end{aligned}}
\]

% manuscript lines / position: 11–13
(car pour deux inv. $\sigma$ et $\tau$,

% manuscript lines / position: 12–13
\[
 \begin{aligned}
 \sigma\tau=\tau\sigma &\iff (\sigma\tau)^2=\deleted{\unsure{\sigma}}\;1\quad\text{i.e.}\\
 &\hphantom{\iff}\sigma\tau\sigma\tau=1
 \end{aligned}\quad\text{)}
\]

% manuscript lines / position: 14–15
En utilisant $1^\circ)$, \deleted{\unsure{ceci}} \inserted{la condition $\sigma_1\sigma_2\sigma_3=1$} s'explicite par la formule

% manuscript lines / position: 16
\[
 2x=T_1(x)+T_2(x)+T_3(x)-\underline{t_2T_{2}(x)}\unsure{))},
\]

% manuscript lines / position: 17–18
en apparence asymétrique, \deleted{mais} elle est équivalente à $\sigma_2\sigma_3\sigma_1=1$ ou $\sigma_3\sigma_1\sigma_2=1$,

% manuscript lines / position: 19–20
\[
 \begin{aligned}
 \text{donc}\qquad 2x&=T_1(x)+T_2(x)+T_3(x)-t_3T_3(x)\\
 2x&=\qquad -\;\;-\;\;-\;\;\text{——}\quad -t_1T_1(x),
 \end{aligned}
\]

% manuscript lines / position: 21
donc elle \uline{implique}\textsuperscript{*} la condition

% manuscript lines / position: 21–24, left, oblique
% slanted
\margin{du moins si $k\longrightarrow A$ injectif --- mais on la \unsure{supposera} vérifiée en tout cas\ldots}

% manuscript lines / position: 22
\[
 \boxed{t_1\mathrm{Tr}_1(x)=t_2\mathrm{Tr}_2(x)=t_3\mathrm{Tr}_3(x)}
\]

% manuscript lines / position: 23–25
(qui ressemble à $6^\circ$ ci-dessus). \unsure{D'on|D'où|Si on} \unsure{désigner|désignant|désigne} par $T_{A/k}(x)$ la valeur commune aux trois mbres, on trouve donc la condition

% manuscript lines / position: 26
\[
 \boxed{2x=T_1(x)+T_2(x)+T_3(x)-T_{A/k}(x).}
\]

% manuscript lines / position: bottom left
% archive sheet number: 113
\scanpage{115}

% manuscript lines / position: 1–3, crossed block
\deleted{Sous réserve que les $\alpha_i:A_i\unsure{\longrightarrow|\mathord{-}} A$ soient bien injectifs, cela montre que chacun des $T_i$ est déterminé par les deux autres, via ``}

% manuscript lines / position: 4, crossed block
\deleted{\[
 \alpha_1T_1(x)=2x+\deleted{\unsure{\alpha_2T_2(x)-\alpha_3T_3}}\inserted{T_{A/k}(x)-\alpha_2T_2(x)-\alpha_3T_3(x)}
\]}

% manuscript lines / position: 5–6, crossed block
\deleted{On s'attend à trouver \deleted{des} formules similaires, multiplicatives}

% manuscript lines / position: 7, crossed block
\deleted{\[
 \boxed{N_{A/k}(x)\,x^2=N_1(x)N_2(x)N_3(x)}
\]}

% manuscript lines / position: 8, crossed block
\deleted{\[
 \text{et}\qquad \deleted{\unsure{N_2}}\,\boxed{n_1(N_1(x))=n_2(N_2(x))=n_3(N_3(x))},
\]}

% manuscript lines / position: 9–11
Mais \deleted{toutes} les \deleted{conditions} \inserted{relations sur la norme} se prouvent aisément, comme \uline{conséquences} de celles sur les traces, i.e. la condition $\sigma_1\sigma_2\sigma_3=1$ sur les $\sigma_i$. En effet, posons

% manuscript lines / position: 12
\[
 x_0=x,\quad x_1=T_1(x),\quad x_2=T_2(x),\quad x_3=T_3(x)
\]

% manuscript lines / position: 13–15
On a alors
\diagram{\(\begin{array}{lll}
 \sigma_1\unsure{:} & x_0\Leftrightarrow x_1 & x_2\leftrightarrow x_3\\
 \sigma_2: & x_0\leftrightarrow x_2 & x_1\unsure{\longleftrightarrow|\longrightarrow}x_3\\
 \sigma_3: & x_0\leftrightarrow x_3, & x_1\unsure{\longleftarrow|\longleftrightarrow}x_2
 \end{array}\)}

% manuscript lines / position: 16–18
\[
 \begin{aligned}
 \text{d'où}\quad &\deleted{N_1(N_1(x))=N_1(x)\overline{N_1(x)}=x\unsure{\sigma_1(x)}(\unsure{x\sigma_1(x)})}\\
 &\deleted{\unsure{N_1(x_1)=x_1\bar x_1}}\\
 n_1(N_1(x))&\overset{!}{=}\underset{\substack{\Vert\\x\deleted{\unsure{T}}\sigma_1(x)}}{\underbrace{N_1(x)}}\,\overline{N_1(x)}
 =(x_0x_1)\underbrace{\sigma_2(x_0x_1)}_{x_2x_3}=x_0x_1x_2x_3
 \end{aligned}
\]

% manuscript lines / position: 19–21
et de $\unsure{=|\text{même}}$ pour $n_2(N_2(x))$, $n_3(N_3(x))$. Donc on a l'égalité des $n_i(N_i(x))$ si $k\unsure{\mathord{-}|\longrightarrow}A$ est \uline{injectif}. Et on aura

% manuscript lines / position: 22–23
\[
 \begin{aligned}
 N_1(x)N_2(x)N_3(x)&=(x_0x_1)(x_0x_2)(x_0x_3)\\
 &=x_0^2\underbrace{(x_0x_1x_2x_3)}_{N_{A/k}(\unsure{x_0|x})}=x^{\unsure{2}}N_{A/k}(x),
 \end{aligned}
\]

% manuscript lines / position: 24
O.K.

% manuscript lines / position: bottom left
% archive sheet number: 114
\scanpage{116}

% manuscript lines / position: 1
Considérons l'algèbre produit tensoriel \unsure{à|:|triple}

% manuscript lines / position: 2
\[
 A_1\otimes A_2\otimes A_3\,;
\]

% manuscript lines / position: 3–5
Via les $\sigma_{A_i}$, le groupe $\mathbb F_2^3$ opère dessus. Il n'y a pas en général d'opération de ce groupe sur $A$, compatible avec

% manuscript lines / position: 6
\[
 \bigotimes A_i\unsure{\longrightarrow|\mathord{-}}A,
\]

% manuscript lines / position: 7–12
\deleted{car} il faudrait qu'il y ait p.ex. un autom. $\tau_1$ de $A$ qui induise \deleted{\unsure{Id|l'}} l'identité sur $A_2,A_3$, et $\sigma_{A_1}$ sur $A_1$; Mais dans les cas les plus simples, $A$ est engendré par $A_2,A_3$, donc $\tau_1$ devrait être l'identité, alors qu'en général $\sigma_{A_1}$ sur $A_1$ ne l'est pas.

% manuscript lines / position: 13
Par contre, soit

% manuscript lines / position: 14–16
\[
 \begin{gathered}
 \unsure{\gamma|\mathcal G|\Gamma}\deleted{\unsure{=}}\subset\mathbb F_2^3\\[-2pt]
 \Vert\\[-2pt]
 \operatorname{Ker}\bigl(\mathbb F_2^3\longrightarrow\mathbb F_2\bigr)\\[-2pt]
 (\deleted{\unsure{x}}\inserted{\unsure{\xi|\varepsilon}}_i)\longmapsto\sum\xi_i
 \end{gathered}
\]

% manuscript lines / position: 17–20
formé des éléments $(\xi_1,\xi_2,\xi_3)$ tels que le nb des composantes $\ne0$ (i.e. $=1$) soit $0$ ou $2$ i.e soit pair. \deleted{\unsure{On|Les}} il y \deleted{\unsure{est}} a trois élmts non nuls de ce groupe, savoir

% manuscript lines / position: 21–23
\[
 \begin{aligned}
 \deleted{(0,0,1)},\qquad &(0,1,1)\\
 &(1,0,1)\\
 &(1,1,0)
 \end{aligned}
\]

% manuscript lines / position: 24
qui correspondent aux autom

% manuscript lines / position: 25–27
\[
 \begin{aligned}
 \mathrm{id}_{A_1}\otimes\sigma_{A_2}\otimes\sigma_{A_3}&\overset{\text{déf}}{=}\tau_1\\
 \sigma_{A_1}\otimes\mathrm{id}_{A_2}\otimes\sigma_{A_3}&\overset{\text{def}}{=}\tau_2\\
 \sigma_{A_1}\otimes\sigma_{A_2}\otimes\mathrm{id}_{A_3}&\overset{\text{def}}{=}\tau_3
 \end{aligned}
\]

% manuscript lines / position: bottom left
% archive sheet number: 115
\scanpage{117}

% manuscript lines / position: 1–2
et \deleted{\unsure{les|ici}} on a une \deleted{\unsure{bonne}} op. de $\unsure{\gamma|\mathcal G|\Gamma}$ sur $A$, telle que

% manuscript lines / position: 2
\[
 \bigotimes A_i\unsure{\longrightarrow|\mathord{-}}A
\]

% manuscript lines / position: 3–5
soit compatible avec des opérations de $\gamma$, savoir celle qui fait opérer \uline{$\tau_i$} par $\sigma_{A/A_i}$.

% manuscript lines / position: 6–10
\deleted{\unsure{Considér|Considérons}} Supposons maintenant que les $A_i$ correspondent \deleted{\unsure{à un|à}} des rev. quadratiques $X_i$ de $X$, donc on a un rev. quadratique $X_1\wedge X_2\wedge X_3$, d'algèbre $B=A_1\unsure{*|\star|\wedge}A_2*A_3$,

% manuscript lines / position: 11–12
\[
 \begin{aligned}
 \text{et on a}\qquad X_1\times X_2\times X_3&\longrightarrow X_1\wedge X_2\wedge X_3,\\
 (x_1,x_2,x_3)&\longrightarrow x_1*x_2*x_3,
 \end{aligned}
\]

% manuscript lines / position: 13
\[
 \text{i.e.}\qquad B\longrightarrow A_1\otimes A_2\otimes A_3,
\]

% manuscript lines / position: 14
et de façon plus précise, on aura

% manuscript lines / position: 15–16
\[
 B\unsure{\longrightarrow|\mathord{-}}(A_1\otimes A_2\otimes A_3)^\gamma
 \qquad\substack{\text{invariants}\\\text{sous }\gamma}
\]

% manuscript lines / position: 17–20
Et \unsure{$=$|de même|même}, si on \unsure{exclut|a exclut} les $x\in X=\operatorname{Spec}(k)$ de \uline{car. 2} tels que deux au moins des $A_i$ soient non étales en ce point, alors l'hom précédent est un \uline{iso}.

% manuscript lines / position: 21–22
On va \uline{regarder} le composé

% manuscript lines / position: 22
\[
 B\unsure{\mathord{-}|\longrightarrow}A_1\otimes A_{\unsure{2|3}}\otimes A_3\unsure{\mathord{-}|\longrightarrow}A,\qquad\deleted{\unsure{\text{i.e.}}}
\]

% manuscript lines / position: bottom left
% archive sheet number: 116
\scanpage{118}

% manuscript lines / position: 1–2
i.e., en termes de $\unsure{\mathscr X|\mathcal X}=\operatorname{Spec} A$, le composé

% manuscript lines / position: 3
\[
 \mathscr X\unsure{\mathord{-}|\longrightarrow} X_1\times_X X_2\times_X X_3\longrightarrow X_1\wedge X_2\wedge X_3,
\]

% manuscript lines / position: 4–7
j'ai envie de prouver que c'est \unsure{un|une} hom \uline{constant} i.e. qu'il se factorise par une \uline{section} de $X_1\wedge X_2\wedge X_3$, ce qui signifie aussi que l'hom composé précédent

% manuscript lines / position: 8–9
\diagram{\begin{tikzcd}[ampersand replacement=\&,column sep=small,row sep=small]
 B \arrow[rr,"f"] \arrow[dr] \& \& A\\
 \& \left(\bigotimes A_i\right)^{\unsure{\gamma|\mathcal G|\Gamma}}\arrow[ur] \&
\end{tikzcd}}

% manuscript lines / position: 10–13
est à valeurs dans $k.1_A$. (Du moins, \unsure{il y équivaut|il y a équivalence|on y aboutira} quand on suppose $k\unsure{\mathord{-}|\longrightarrow}A$ injectif --- c'est ce qu'on va supposer maintenant, pour simplifier.)

% manuscript lines / position: 14
Mais on a évidemment

% manuscript lines / position: 15
\[
 f(B)\subset A^\gamma,
\]

% manuscript lines / position: 16
et la condition est vérifiée si on suppose

% manuscript lines / position: 17
\[
 \boxed{A^\gamma=k.\unsure{1|\bar 1}_A}.
\]

% manuscript lines / position: 18–21
(C'est là une condition pas tjrs vérifiée, \unsure{$=$|même} dans des cas intéressants, à cause de la car. 2. Mais si 2 est inversible \deleted{\unsure{,}} dans $k$, alors tjrs vrai \unsure{compte tenu|au terme} des conditions \unsure{précdtes|précédentes}).

% manuscript lines / position: bottom left
% archive sheet number: 117
\scanpage{119}

% manuscript lines / position: 1–2
En effet, si $x\in A$, $\sigma_i(x)=x$ signifie $\deleted{\unsure{x=}}T_i(x)-x=x$ i.e.

% manuscript lines / position: 3
\[
 2x=T_i(x),
\]

% manuscript lines / position: 4
et si c'est vrai pour $i=1,\unsure{2|1},3$ on trouve

% manuscript lines / position: 5–7
\[
 6x=\underbrace{T_1(x)+T_2(x)+T_3(x)}_{\substack{\Vert\ \text{cf plus haut}\\ 2x+T_{A/k}\,\unsure{x|(x)}}}
\]

% manuscript lines / position: 8
\[
 \text{d'où}\qquad\boxed{4x=\deleted{\unsure{k}}T_{A/k}(x)}
\]

% manuscript lines / position: 9–10
\[
 \begin{aligned}
 &\text{Ce qui implique, si 2 inv}\\
 &\qquad x=\frac14 T_{A/k}(x)\in k\qquad\text{OK.}
 \end{aligned}
\]

% manuscript lines / position: 11–12, crossed block
\deleted{Dans le cas général, on va supposer cette condition $A^{\unsure{\gamma|\mathcal G|\Gamma}}=k.1_A$.}

% manuscript lines / position: 13–17
\uline{Proposition} Supposons les $A_i$ des alg. \unsure{quadratiques|quadratique} \uline{étales}, que $k\longrightarrow A$ soit injectif, que les $\unsure{i_i|\iota_i|\alpha_i}(A_i)\subset A$ engendrent l'algèbre $A$ i.e. $\bigotimes A_i\unsure{\mathord{-}|\longrightarrow}A$ épi, (donc $\unsure{\mathscr X|\mathcal X}=\operatorname{Spec} A\longrightarrow X_1\times_k X_2\times_k X_3$ est une immersion fermée)\inserted{, enfin $A^\gamma=k.1_A$.} Alors \deleted{la situation} \inserted{les \unsure{hom|flèches}} \deleted{\unsure{la situation}} \deleted{se déduit de la situation \unsure{standard}}

% manuscript lines / position: 18
\[
 \mathscr X\longrightarrow X_1\times X_2,\quad \mathscr X\longrightarrow X_2\times X_3,\quad \mathscr X\longrightarrow X_3\times X_1
\]

% manuscript lines / position: 19–23
sont des iso. (donc $\mathscr X$ est \inserted{fini} \uline{étale} sur $X=\operatorname{Spec} k$). Et la situation se déduit \uline{par twist} de celle où 3 ens. à deux élmts \deleted{\unsure{$\xi$}}, $I_1,I_2,I_3$, munis d'une trivialisation $\varepsilon$ de $I_1\wedge I_2\wedge I_3$, et \unsure{du|de} $J\subset I_1\times I_2\times I_3$,

% manuscript lines / position: 24, at foot
\[
 J=\{(\alpha_1,\alpha_2,\alpha_3)\in I_1\times I_2\times I_3\mid \alpha_1\wedge\alpha_2\wedge\alpha_3=\varepsilon\}.
\]

% manuscript lines / position: bottom left
% archive sheet number: 118
\scanpage{120}

% manuscript lines / position: 1–5
Inversement, \inserted{dans} une telle situation twistée $A$ est une algèbre étale quadratique sur chacun des $A_i$, donc on a une situation biquadratique \inserted{étale}, et de plus les hypothèses du début (en particulier $A^{\unsure{\gamma|\sigma|\Gamma}}=k.1_A$) sont satisfaites.

% manuscript lines / position: 6–10
\deleted{Il faut} Seule la première partie, directe, demande un peu de soin. OPS $X_1,X_2,X_3$ rev. triviaux, $I_{1X},I_{2X},I_{3X}$, et \deleted{la section} \inserted{la section construite plus haut de $X_1\wedge X_2\wedge X_3=(I_1\wedge I_2\wedge I_3)_X$} constante. \deleted{\unsure{on}} On aura alors

% manuscript lines / position: 11
\[
 \unsure{\mathscr X|\mathcal X}\subset J_X
\]

% manuscript lines / position: 12
et tout résulte des lemmes suivants

% manuscript lines / position: 13–15
\uline{Lemme} Soit $\unsure{\gamma|\mathcal G|\Gamma}$ un groupe, opérant de façon transitive sur un ens. $J$, et considérons un sous-schéma $\mathscr X$ de $X_J$ \inserted{(où $X=\operatorname{Spec} k$)} qui soit tel que

% manuscript lines / position: 16
a) $\mathscr X$ est stable sous $\gamma$, et

% manuscript lines / position: 17–18
b) le composé $k\longrightarrow\Gamma(X_J,\unsure{\underline O|\mathcal O}_{X_J})\;(\simeq k^J)$ est injectif.

% manuscript lines / position: 19–20
Alors $\mathscr X=X_J$, donc $\mathscr X$ est étale sur $X$ (et fini étale si $J$ est fini).

% manuscript lines / position: 21–24
En effet, soit $j_0\in J$, considérons la composante $X_{\unsure{0|j_0|j}}$ de $X_J$, $X_j\xrightarrow{\sim}X$, et $\mathscr X_j=X\cap\mathscr X$, il correspond à un sous-schéma fermé de $X$ donc à un idéal $\unsure{\mathfrak a|\mathcal A|\mathfrak J_j}$ de $k$. A cause de l'\unsure{homog|action} de

% manuscript lines / position: bottom left
% archive sheet number: 119
\scanpage{121}

% manuscript lines / position: 1–5
$\unsure{\mathscr X|\mathcal X}$ sous $\unsure{\gamma|\mathcal G|\Gamma}$, et de la transitivité de $\gamma$ sur \unsure{$I$|$J$}, on voit que cet idéal $\unsure{\mathfrak a|\mathcal A|\mathfrak J}_j$ ne dépend pas de $j$. D'autre part, l'intersection des $\mathfrak a_j$ est évidemment le noyau de $\deleted{\unsure{\mathfrak a}}k\longrightarrow\Gamma(X_J,\unsure{\underline O|\mathcal O}_{X_J})\simeq k^J$. Comme celui-ci est nul, on a $\mathfrak a_j=0$ donc

% manuscript lines / position: 6
\[
 \mathscr X_j=X_j\quad\text{donc}\quad\mathscr X=X_J\qquad\text{O.K.}
\]

% manuscript lines / position: 7–13
Ce résultat vaut sans restriction de caractéristique, mais il concerne la situation standard stricte, \uline{très} restrictive. J'ai envie de vérifier que dès qu'on a trois rev. quadratiques $X_1,X_2,X_3$ de $X$, et une section $\varepsilon$ de $X_1\wedge X_2\wedge X_3$, sans condition étale, on en déduit sur

% manuscript lines / position: 14
\[
 X=\pi^{-1}(\varepsilon)\qquad
 \deleted{\unsure{\pi}}\; X_1\times_X X_2\times_X X_3\xrightarrow{\pi}X_1\wedge X_2\wedge X_3
\]

% manuscript lines / position: 15–17
une structure biquadratique p.r. aux $\unsure{X_i|A_i}$. Ceci doit généraliser la structure biquadratique can. sur un $X_1\times X_2$, par les applications

% manuscript lines / position: 18–20
\diagram{\begin{tikzcd}[ampersand replacement=\&,row sep=small]
 \& X_1\\
 X_1\times X_2 \arrow[ur] \arrow[r] \arrow[dr] \& X_2\\
 \& X_3=X_1\wedge X_2,
\end{tikzcd}}

% manuscript lines / position: 21
comme on va le voir \unsure{tantôt|bientôt|de suite}.

% manuscript lines / position: bottom left
% archive sheet number: 120
\scanpage{122}

% manuscript lines / position: 1–3
On va regarder, \deleted{pour} $X_1,X_2,X_3$ fixés, la situation universelle, où la \inserted{(nouvelle)} base sera $\unsure{\mathscr X|\mathcal X|\mathfrak X}_0=X_1\wedge X_2\wedge X_3$, \deleted{où} \unsure{muni} des \unsure{rev|revêtement}\deleted{\unsure{. ét}}

% manuscript lines / position: 4
\[
 X_{1\mathscr X_0},\quad X_{2\mathscr X_0},\quad X_{3\mathscr X_0}
 \qquad (X_{i\mathscr X_0}=X_i\times_X\mathscr X_{\unsure{0|1}})
\]

% manuscript lines / position: 5
et où $\mathscr X$ est $X_1\times X_2\times X_3$, avec

% manuscript lines / position: 6–8
\diagram{\begin{tikzcd}[ampersand replacement=\&,column sep=large,row sep=small]
 \& X_{1\mathscr X_0} \arrow[dr] \&\\
 X_1\times X_2\times X_3 \arrow[ur,"{(\mathrm{pr}_1,\wedge)}"] \arrow[r,"{(\mathrm{pr}_2,\wedge)}"] \arrow[dr,"{(\mathrm{pr}_3,\wedge)}"'] \& X_{2\mathscr X_0} \arrow[r] \& \mathscr X_0\mathord{-}\\
 \& X_{3\mathscr X_0} \arrow[ur] \&
\end{tikzcd}}

% manuscript lines / position: 9–11
où \deleted{\unsure{$X_1$}} $X_1\times X_2\times X_3\unsure{\longrightarrow|\mathord{-}}\mathscr X_0=X_1\wedge X_2\wedge X_3$ est le morphisme

% manuscript lines / position: 12
\[
 x_1,x_2,x_3\longmapsto x_1\wedge\unsure{\alpha_2|x_2}\wedge x_3
\]

% manuscript lines / position: 13–15
\unsure{Une} difficulté, c'est que ce morphisme n'est pas néc fini loc libre, ni $=$ plat (mais il est bien sûr \unsure{fini}).

% manuscript lines / position: 16–19
\diagram{\begin{tikzcd}[ampersand replacement=\&,column sep=huge,row sep=large]
 B=A_1\otimes_{O}A_2\otimes_O A_3
 \arrow[r,bend left=55,"{\mathrm{Tr}_1}"]
 \arrow[r,bend left=35,"{N_1\,?}"]
 \arrow[r,shift left=5pt]
 \arrow[r,shift left=2pt,"{\deleted{\unsure{\mathrm{Tr}_2,N_2}}}" description]
 \& A_1\otimes_O B_0=B_1\quad\unsure{,} B_0=A_1*A_2*A_3 \arrow[l,shift left=2pt,"{2(!)}"]\\
 B_0 \arrow[u,"{4(!)}"'] \arrow[u,phantom,"{\text{(\unsure{pas néc. inj.})}}" {xshift=2em}] \arrow[ur,"{\unsure{2|\varepsilon 2}}"] \&\\
 O\arrow[u,"2"']\arrow[uu,bend left=48,no head,"8"] \&
\end{tikzcd}}

% manuscript lines / position: 20
$\mathrm{Tr}_1$ doit être $B_0$-linéaire, $N_1$ $B_0$-quadratique\ldots

% manuscript lines / position: bottom left
% archive sheet number: 121
\scanpage{124}

% manuscript lines / position: 1–4
\diagram{\begin{tikzcd}[ampersand replacement=\&,column sep=small]
 0\arrow[r]\&L\arrow[r]\arrow[d,"\alpha"']\&E\arrow[l,bend right=35,"\pi"']\arrow[r,no head]\arrow[d,no head,"{\beta\,\unsure{\uparrow|\downarrow}}"]\&M\arrow[l,bend right=35,"{\unsure{\varpi|\bar\omega}}"']\arrow[r,no head]\arrow[d,"\gamma",no head]\&0\\
 0\arrow[r,no head]\&L'\arrow[r,no head]\&E'\arrow[l,bend left=35,"\pi'"]\arrow[r,no head]\&M'\arrow[l,bend left=35,"\varpi'"]\arrow[r,no head]\&0
\end{tikzcd}}

% manuscript lines / position: 5
\[
 \alpha\pi(x)\unsure{\equiv|=}\pi'\beta(x)\overset{?}{\Longleftrightarrow}\beta\varpi(\unsure{\eta|y})=\varpi'\gamma(y)\;?
\]

% manuscript lines / position: 6–8
\[
 \begin{aligned}
 \underbrace{\deleted{\beta}\alpha\pi(x)=\pi'\beta(x)}&\Longleftrightarrow\beta\varpi(x)\equiv\\
 \beta(\unsure{\mathrm{id}}-\varpi)&=(\deleted{\unsure{\pi}}\unsure{\mathrm{id}}-\varpi')\beta\\
 &\deleted{\unsure{-\beta\varpi=-\varpi'\beta}}
 \end{aligned}
\]

% manuscript lines / position: 9–10
\diagram{\begin{tikzcd}[ampersand replacement=\&,column sep=small]
 E_1\arrow[d]\& E_2\arrow[d]\&\Rightarrow\& E_1\wedge E_2\arrow[r]\&E'_{\unsure{1|i}}\\
 E'_1\&E'_2\&\&\&
\end{tikzcd}}

% manuscript lines / position: 11–13
\diagram{\scriptsize\begin{tikzcd}[ampersand replacement=\&,column sep=small,row sep=small]
 \& E_1\otimes E_2\arrow[r]\& E'_1\otimes\deleted{\unsure{E}}E'_{\unsure{2}}\arrow[dr]\arrow[dr,shift left=2pt,shorten <=7pt,shorten >=5pt]\&\\
 {E_1\otimes L_2,\ L_1\otimes E_2}\arrow[ur]\arrow[dr]\&\&\& E'_1\wedge E'_2\\
 \&L_1\otimes L_2\arrow[r]\&L'_1\otimes L'_2\arrow[ur,"{\text{inj can}}"']\&
\end{tikzcd}}

% manuscript lines / position: 14–17
\diagram{\scriptsize\begin{tikzcd}[ampersand replacement=\&,column sep=tiny,row sep=large]
 \unsure{x_1}\otimes\unsure{u_2|x_2}\arrow[r,mapsto]\arrow[dr,bend right=45]\&
 \beta_1(x_1)\otimes\alpha_2(u_2)\arrow[r]\&
 \begin{gathered}\beta_1(x_1)*\alpha_2(u_2)\\
 =\pi'_1\deleted{\unsure{\pi}}\beta_1(x_1)\otimes\deleted{\unsure{\pi_2|\gamma|\beta}}\alpha_2(\unsure{\mathrm{id}(u_2))|\beta_2(x_2))})\end{gathered}\\
 \&\pi_1(x_1)\otimes u_2\arrow[r]\&
 \underbrace{\alpha_1\pi_1(x_1)}_{\pi_1\beta_1(x_1)}\otimes\alpha_2(x_2)
\end{tikzcd}}

% manuscript lines / position: 18
\[
 E_1\wedge
\]

% manuscript lines / position: bottom left
% archive sheet number: 123
\scanpage{125}

% manuscript lines / position: top centre
\mspage{465}

% manuscript lines / position: 1–10
Il me semble que notre époque se caractérise comme celle d'une exacerbati\unsure{on} à outrance de cette dégradation culture\deleted{\unsure{bt|b}}\unsure{ll}e. Parmi les derniers actes de cette histoire, il y a ceux, intimement solidaires, de la ``course à l'espace'' entre le\unsure{s} deux superpuissances antagonistes (imbues de valeurs essentiellement identiques) et de la course aux armements (nucléaires notamment). Comme acte ultime et dénou\unsure{e}ment probable de cette évolution forcenée dans les surenchères d'un certain type de ``force'' ou de ``pouvoir'', on peut prévoir dès à présent quelque holocauste nucléaire (ou autre, il y a l'embarras du choix...) à l'échelle planétaire. Il aur\unsure{a} peut-être ce mérite de résoudre tous les problèmes d'un seul coup et une fois po\unsure{ur} toutes...

% manuscript lines / position: 11–25
Mon propos ici n'est pas pourtant de brosser un alléchant tableau de ``fin du monde'' (on ne m'a pas attendu pour cela), et encore moins de partir en guerre contre le ``muscle'', ou contre ``le cerveau'' (alias l'``esprit\unsure{''}). Je sais bien que même mes ``tripes'' n'auraient rien à y gagner! Je tiens à mes muscles et à mon cerveau, qui me sont bien utiles on s'en doute, comme je tiens aussi à mes ``tripes'', qui ne le sont pas moins. Plutôt, il me semble utile de dire ici en quelqu\unsure{es} mots (si faire se peut) comment s'est joué dans ma propre personne ce conflit pr\unsure{o}fond, véhiculé par la culture environnante, entre ces deux types de valeurs. En termes plus terre à terre, il s'agit aussi de l'historique de mes attitudes (d'acceptation voire d'exaltation, ou de rejet) de deux \mbox{a s p e c t s} ou \mbox{f a c e s} également réels et tangibles de ma personne, inséparables et complémentaires par nature\handwritten{\unsure{,}} et nullement antagonistes par eux-mêmes. Je pourrais les appeler ``\mbox{l' h o m m e}'' et ``\mbox{l a\quad f e m m e}'' en moi, ou aussi (pour prendre des appe\handwritten{\unsure{ll}}ations moins ``chargées'', et qui pour cela offrent moins de risque d'induire en erreur), le ``\mbox{y a n g}'' et le ``\mbox{y i n}''.

% manuscript lines / position: 26–35
Il semblerait que chez la plupart des personnes, les ``jeux sont faits'' dès la petite enfance, où se mettent en place les mécanismes essentiels qui, pendant la vie entière, vont dominer en silence, avec une efficacité d'automate parfaitement au point, attitudes et comportements. Au coeur de\unsure{.} ces mécanismes sont ceux d'affirmation ou de rejet de tels et tels traits en nous, ou de telles pulsions profondes, à ``signature'' soit yang soit yin, ou de tels \unsure{et tels} ``paquets'' de traits et de pulsions à signature donnée, voire même du paquet ``yang'' ou du paquet ``yin'' tout entier. Ce sont ces mécanismes qui, dans une très large mesure, déterminent tous les autres mécanismes de choix (affirmation ou rejet) structura\unsure{nt} notre ``moi''.

% manuscript lines / position: bottom left
% archive sheet number: 124
\scanpage{126}

% manuscript lines / position: 1–4
\diagram{\scriptsize\begin{tikzcd}[ampersand replacement=\&,column sep=small,row sep=large]
 {\unsure{\mathcal E|\mathscr E}_{(p_1,q_1)(p_2,q_2)}=(M_1\otimes M_2)\oplus(L_1\otimes L_2)}
 \arrow[rr,"{T_U}"] \arrow[dd,"{T_{-\gamma b}}"'] \& \& {\mathcal E_{(p'_1,q'_1)(p'_2,q'_2)}}\arrow[dd,no head,"{T_{-\gamma b'}}"]\\
 \& |[draw,ellipse]| \text{\unsure{commu|commute}!} \&\\
 \mathcal E \arrow[rr,"{T_{-U}}"] \& \& \mathcal E
\end{tikzcd}}

% manuscript lines / position: 5–6
\[
 \begin{aligned}
 b'&=b\deleted{\unsure{-}}+\chi U\\
 U-\gamma b'&=U-\gamma b\deleted{\unsure{=}}\unsure{-}\gamma\chi U=-U-\gamma b
 \end{aligned}
\]

% manuscript lines / position: 7
d'où un iso \inserted{can} d'ext. $\chi$-scindées

% manuscript lines / position: 8–9
\[
 \begin{aligned}
 E_1\vee E_2&\simeq E_1\wedge E_2\\
 E'_1\vee E'_2&\simeq E'_1\wedge E'_2
 \end{aligned}
\]

% manuscript lines / position: 10–13
Je dis qu'il est fonctoriel (pour les \uuline{iso}): c'est clair par construction. Mais il y a aussi une notion d'\uuline{hom} d'ext. $\chi$-scindées. Ça marche aussi

% manuscript lines / position: 14–17, lower left
\diagram{\scriptsize\begin{tikzcd}[ampersand replacement=\&,column sep=small,row sep=small]
 \& L_i\arrow[r,no head]\arrow[dl,equal]\arrow[dd,no head]\&
 E'_i\arrow[r,"{\deleted{\unsure{M'_i}}}"]\arrow[dl,dashed]\arrow[dd]\& {M'_i\arrow[dl]\arrow[dd,equal]}\\
 L_i\arrow[r,no head]\arrow[dd]\& E_i\arrow[r]\arrow[dd]\& M_i\arrow[dd,equal]\&\\
 \& L'_i\arrow[r,no head]\arrow[dl,equal]\&\unsure{E'''_i|E''_i}\arrow[dl,dashed]\arrow[r,no head]\& M'_i\arrow[dl]\\
 L'_i\arrow[r]\&E''_i\arrow[r]\&M_i\&
\end{tikzcd}}

% manuscript lines / position: 14–17, lower right
\diagram{\scriptsize\begin{tikzcd}[ampersand replacement=\&,column sep=small,row sep=small]
 \&\&\&M'_1\otimes M'_{2}\arrow[dl]\\
 L_1\otimes L_2\arrow[d]\arrow[r]\& E\arrow[r,two heads]\&M_1\otimes M_2\&\\
 L'_1\otimes L'_2\&\&\&
\end{tikzcd}}

% manuscript lines / position: bottom left
% archive sheet number: 125
\scanpage{127}

% manuscript lines / position: top centre
\mspage{466}

% manuscript lines / position: 1–14
Pour des raisons qui restent encore mystérieuses pour moi, dans mon propre cas l'histoire des relations (tant conscientes qu'inconscientes) entre le moi (``le patron'')\handwritten{,} et ``le mâle'' et ``le féminin'' en ma personne (aussi bien dans le ``patron'' lui-même que dans l'``ouvrier'', qui l'un et l'autre sont tributaires du double aspect yin-yang de toutes choses) - cette histoire a été plus mouvementée qu'à l'accoutumé\handwritten{e}. J'y distingue trois périodes. La dernière rejoint dans un certain sens la première, qui s'étend sur les cinq premières années de mon enfance. Cette troisième période, que je peux appeler celle de la \mbox{m a t u r i t é} , peut être vue comme une sorte de ``retour'' à cette enfance, ou comme de progressives retrouvailles avec l' ``\mbox{é t a t\quad d'\quad e n f a n c e}'', avec l'harmonie des épousailles sans histoires du ``yin'' et du ``yang'' en mon être. Ces retrouvailles ont commencé au mois de juillet 1976, à l'âge de quarante huit ans - l'année même où j'ai fait la découverte (trois mois plus tard) d'un pouvoir jusque là ignoré en moi, le pouvoir de méditation (*).

% manuscript lines / position: 15–32
Les valeurs dominantes dans la personne de chacun de mes parents, tant ma mère que mon père, étaient des valeurs yang: volonté, intelligence (au sens: pui\unsure{s}sance intellectuelle), contrôle de soi, ascendant sur autrui, intransigeance, ``Konsequenz'' (qui signifie, en allemand, cohérence extrême dans (ou avec) ses options, idéologiques notamment), ``idéalisme'' au niveau politique comme pratique... Chez ma mère, cette valorisation a pris dès son jeune âge une force exacerbée, c'était le revers d'une véritable haine qu'elle avait développée vis à vis de ``la femme'' en elle (et à partir de là, vis-à-vis du féminin en général\unsure{:}\handwritten{)}. Cette haine en elle a fini par prendre une véhémence et une force d'autant plus destructice, qu'elle est restée entièrement occulté\handwritten{e} sa vie durant.\handwritten{(}Moi-même ai fini par découvrir ces choses il y a cinq ans seulement, trois ans après que la méditation apparaisse dans ma vie.) Dans un tel contexte parental, c'est un mystère (et pourtant un fait qui ne fait aucun doute pour moi) que j'aie pu pleinement m'épanouir pendant les cinq premières années de mon enfance - jusqu'a\unsure{u} moment de l'arrachement au milieu parental et de la destruction de ma famille d'origine (formée de mes parents, de ma soeur plus agée, et de moi), de par la volonté de ma mère et à la faveur (si on peut dire) des évènements politiques de l'année 1933.

% manuscript lines / position: footnote, beneath rule
\noindent\rule{.24\linewidth}{.35pt}\par
(*) Voir les deux sections ``Désir et méditation'' et ``L'émerveillement'', n\textsuperscript{o}s 36 et 37.

% manuscript lines / position: bottom left
% archive sheet number: 126
\scanpage{128}

% manuscript lines / position: 1–2
\[
 \begin{aligned}
 \tilde b_i&=\varpi_i-\chi q_i,\qquad b_i+\tilde b_i=\chi\mathrm{id}-\chi\mathrm{id}=0\\
 &=\underline{\underline{-b_i}}
 \end{aligned}
\]

% manuscript lines / position: 3–4
\[
 \begin{aligned}
 \tilde U_i\ \text{tel que}\quad q'_i&=\underbrace{q_i}_{\mathrm{id}-p_i}-\tilde U_i
 =\deleted{\unsure{\chi}}\mathrm{id}-\underbrace{(p_i+\tilde U_i)}_{p'_i}=\unsure{\cdot}\\
 \tilde U_i&=-U_i
 \end{aligned}
\]

% manuscript lines / position: 5–6
\[
 \begin{aligned}
 \deleted{\unsure{U}}\tilde U&=\tilde b_1\otimes\tilde U_2+\tilde U_1\otimes\tilde b_2+\chi\tilde U_1\otimes\tilde U_2\\
 &=b_1\otimes U_2+U_1\otimes b_2+\chi U_1\otimes U_2=U
 \end{aligned}
\]

% manuscript lines / position: 5–6, right of vertical separator
\margin{Mais Attention, $\tilde T_{\tilde U}=\deleted{\unsure{+}}T_{-\tilde U}=-T_{-U}=T_U^{-1}$}

% manuscript lines / position: 7
\[
 \tilde b=\tilde b_1\tilde b_2=b_1b_2
\]

% manuscript lines / position: 8–11, left
\diagram{\scriptsize\begin{tikzcd}[ampersand replacement=\&,row sep=large]
 {\unsure{\mathcal E|\mathscr E}_{(p_1,q_1)(p_2,q_2)}\deleted{\unsure{\simeq}}=M_1\otimes M_2\unsure{+|\oplus}L_1\otimes L_2}
 \arrow[d,"{T_{v_{(p_1,q_1)(p_2,q_2)}}}","{\sim}"']\\
 {\mathcal E_{(\tilde p_1,\tilde q_1)(\tilde p_2,\tilde q_2)}=M_1\otimes M_2\unsure{+|\oplus}L_1\otimes L_2}
\end{tikzcd}}

% manuscript lines / position: 8–9, right of upper object
$\chi$-scindage défini \inserted{(conv. \unsure{g.|g})} par

% manuscript lines / position: 9, upper right
\[
 M_1\otimes M_2\xrightarrow{b=b_1b_2}L_1\otimes L_2
\]

% manuscript lines / position: 10–11, right of lower object
$\chi$ sc. défini (conv. \unsure{g.}) par $-b$

% manuscript lines / position: 12–14
$T_v$ \inserted{$(=(y,z)\longmapsto y,z\deleted{\unsure{-}}\unsure{+}v(y))$} \deleted{\unsure{tr}} transforme le $\chi$-scindage défini par $b$ en le $\unsure{\phi|\chi}$-scindage défini \deleted{par} $b+\chi v$.

% manuscript lines / position: 15
\[
 \deleted{\unsure{\text{Si}}}\ \text{On veut que}\quad b+\chi v=\underset{\substack{\Vert\\-b}}{\tilde b}\deleted{\unsure{=}}
 \quad\text{i.e.}\quad v\overset{!}{=}-\frac2{\chi}b
\]

% manuscript lines / position: 16–17
Supposons que $\exists\,\gamma\in k$, tq. $\gamma\chi=2$, on peut prendre \deleted{\unsure{$\xi$}}

% manuscript lines / position: 17
\[
 \boxed{v=-\gamma b}
\]

% manuscript lines / position: bottom left
% archive sheet number: 127
\scanpage{129}

% manuscript lines / position: top centre
\mspage{467}

% manuscript lines / position: 1–12
(\textsuperscript{106}1) (3 octobre) Ni moi, ni Deligne n'avons jamais eu le moindre doute que les conjectures de Weil puissent ne pas être valables, et je ne me rappelle pas avoi\unsure{r} entendu quiconque exprimer de tels doutes. Qualifier le ``résultat'' (i.e. la demonstration de ces conjectures) comme ``surprenant'', témoigne encore du propos délibéré d'épater la galerie. D'ailleurs à aucun moment depuis l'introduction de la ``topologie'' et de la cohomologie étales, je n'ai eu le sentiment que ces conjectures étaient hors d'atteinte, mais plutôt \handwritten{(}à partir de 1963\handwritten{)} qu'elles ne manqueraient pas d'être démontrées dans les toutes prochaines années. Au moment de mon départ, en 1970, je n'avais guère de doute que Deligne, qui était le mie\unsure{ux} placé de tous pour cela, ne tarderait pas à les prouver (ce qu'il n'a pas manqué de faire), en même temps que les``conjectures standard sur les cycles algébriques\unsure{'',|''|,} plus fortes (qu'il s'est par contre attaché à discréditer).

% manuscript lines / position: 13–32
C'est d'ailleurs avec raison que Deligne émet des réserves sur la validité de ces dernières conjectures, dont je ne suis pas plus convaincu que lui. Mais la portée d'une conjecture ne dépend pas du fait si elle finira par se révéler vrai\handwritten{e}, ou fausse, pas plus que son caractère de soi-disante ``difficulté'', qui la rendrait``hors de portée'' - caractère entièrement subjectif, lui. Elle dépend un\unsure{i}quement du fait si la \mbox{q u e s t i o n} sur laquelle la conjecture met le doi\unsure{gt} (et qui n'avait pas été perçue, avant qu'elle ne soit posée) - si cette questio\unsure{n} touche à quelque chose de vraiment essentiel pour notre connaissance des choses. Or il est evident (pour moi tout au moins!) qu'il ne saurait être question d'av\unsure{oir} une bonne compréhension des cycles algébriques, ni des propriétés dites ``arithm\unsure{é}tiques'' de la cohomologie des variétés algébriques (ou encore, de la ``géométrie des motifs''), aussi longtemps que la question de la validité de ces conjectures n'est pas réglée. Aujourd'hui encore comme lors du Congrès de Bombay en 1968, je considère cette question, avec celle de la résolution des singularités, comm\unsure{e} l'une des deux questions les plus fondamentales qui se posent en géométrie algébrique. Je sens bien la portée de l'une et de l'autre! Cette fécondité potentie\unsure{l}le ne pourra manquer de se manifester, dès l'instant où on ne se bornera plus à contourner cahin-caha une conjecture décrétée ``trop difficile'', et où quelqu'un prendra la peine enfin de retrousser ses manches et de s'y coltiner!

% manuscript lines / position: archive lower left in scan orientation
% archive sheet number: 128
